\chapter{سربازی اجباری}
\label{ch:military-service}

\section{مقدمه}
متاسفانه در صورتی که پسر باشید در مسیر اپلای تحصیلی و خروج از کشور با مراحل بیشتری مواجه می‌شوید که طی کردن برخی از این مراحل
نیاز به دانستن قوانین و مقررات سربازی و نظام وظیفه دارد. لازم به ذکر است که قوانین نظام وظیفه و سربازی ثابت نیست و بهتر است 
همیشه برای کسب اطمینان از صحت محاسباتتان از دوستان، دانشگاه و \href{https://vazifeh.police.ir/}{\textbf{وبسایت سازمان نظام وظیفه}}
 کسب تکلیف کنید. 

از آنجایی که ذکر کردن کل قوانین و جزئیات قوانین سربازی بی‌ربط به این راهنما و خارج از حوصله آن است، این فصل به صورت بخش بخش و برای 
قوانین مربوط به مهاجرت تحصیلی ذکر می‌شود و هر بخش ممکن است دارای تعدادی سوال متداول باشد. این سوال و جواب‌ها به همراه محتوای موجود 
حاصل ترکیب تجربه شخصی و محتوای استخراج شده از \href{https://t.me/joinchat/BlviDUvKluDK9m3p7OpIxQ}{گروه تلگرامی نظام وظیفه} است و به 
همین جهت از تمام کسانی که در طول سالیان جواب سوالات افراد را داده‌اند تشکر می‌کنیم. 

\section{قوانین و فرایند‌ها}


\subsection{فرایند گرفتن پاسپورت}
\subsubsection{اگر دانشجو هستید}
مراحل ثبت درخواست خروج و اخذ گذرنامه از طریق سامانه‌های دولتی و پلیس +۱۰

\begin{enumerate}[leftmargin=.5in]
    \item افتتاح حساب در بانک سپه\footnote{نیاز نیست وثیقه حتماً از مبدأ بانک سپه واریز شود، اما برای برگشت وثیقه و در مرحله ثبت درخواست به شماره حساب این بانک نیاز دارید.}
    
    \item ثبت درخواست خروج در سامانه سخا، دقت کنید که درخواست خروج دارای \textbf{مقصد}، \textbf{تاریخ شروع} و \textbf{تاریخ پایان} مشخص است و نمی‌توانید به صورت دیگری از آن استفاده کنید
    
    \item تأیید درخواست توسط دانشگاه و سازمان نظام وظیفه، برای تایید درخواست توسط دانشگاه معمولا باید یک فرم درخواست پر کنید
    
    \item پرداخت مبلغ وثیقه پس از تأیید نهایی
    
    \item مراجعه به پلیس +۱۰ جهت تکمیل فرم اخذ پاسپورت و عکس‌برداری. بازه خروج را حداقل ۱۰ تا ۱۵ روز در نظر بگیرید.
    
    \item دریافت پاسپورت به‌صورت پستی و سپس مراجعه حضوری به نظام وظیفه برای آزادسازی وثیقه در صورتی که از کشور خارج نشده‌اید
\end{enumerate}


\begin{tipbox}
    \begin{enumerate}
        \item در برخی بازه‌های زمانی مثل اربعین امکان اخذ پاسپورت بدون وثیقه و در شرایط ساده‌تر وجود دارد اما باید توجه کنید که پاسپورت زیارتی نگیرید.
        \item دقت کنید که برخی دانشگاه‌ها در ترم های آخر، یا در شرایطی که دفاع نکرده باشید در دادن اجازه خروج سخت‌گیری می‌کنند
        \item زمان اتمام مجوز خروج با مراجعه حضوری به نظام وظیفه قابل تمدید است و نیاز به موافقت مجدد دانشگاه ندارد
    \end{enumerate}
\end{tipbox}

\subsubsection{اگر فارغ‌التحصیل هستید}
تکمیل شود.


\subsection{گرفتن برگه سبز/برگه اعزام}
برای گرفتن ریزنمرات رسمی (بعد از فارغ‌التحصیلی) و گرفتن دانشنامه باید تکلیف سربازی اجباری مشخص شود. مشخص کردن تکلیف به معنی گرفتن معافیت دائم یا 
موقت، اعزام به خدمت، مشخص کردن تاریخ اعزام به خدمت یا اشتغال به تحصیل در مقطع بالاتر است. راه‌های گرفتن اشتغال به تحصیل در مقطع بالاتر و مشخص کردن تاریخ
اعزام به‌وسیله برگه‌سبز جزو راه‌های پر استفاده فارغ‌التحصیلان است که در این بخش نحوه گرفتن برگه‌سباز توضیح داده می‌شود.

فرآیند دریافت برگه اعزام به خدمت پس از فارغ‌التحصیلی یا انصراف.

\begin{enumerate}
    \item فارغ‌التحصیلی و اتمام معافیت تحصیلی
    \par\smallskip
    {\small
    پیش از گرفتن برگه‌سبز باید وضعیت معافیت تحصیلی شما در سامانه سخا به عنوان \textbf{مشمول} ثبت شده باشد.
    این اتفاق پس از صدور نامه معرفی به نظام رخ می‌دهد که گاهی اوقات چند روز طول می‌کشد.
    \par}

    \item مراجعه به پلیس +۱۰ با مدارک شناسایی و نامه فراغت از تحصیل یا انصراف.
    \par\smallskip
    {\small
    می‌توان آخرین روز فرجه را به عنوان تاریخ اعزام درخواست داد. می‌توانید پیش از مراجعه به دفتر پلیس+۱۰ این تاریخ 
    را حساب کنید. 
    \par}

    \item دریافت فرم‌های واکسیناسیون و معاینه پزشکی
    
    \item انجام معاینات و تزریق واکسن در مراکز تعیین‌شده
    \par\smallskip
    {\small
    در مرحله معاینه پزشکی، یک فرم درباره سوابق بیماری خود پر خواهید کرد، توجه کنید که در صورتی که سابقه بیماری اعلام کنید یا 
    پزشک تشخیص بیماری اعم از جسمی یا روحی بدهد و آن  را در سامانه ثبت کند به کمیسیون پزشکی فرستاده خواهید شد و در صورت 
    نگرفتن معافیت دائم، نظام وظیفه اولین تاریخ ممکن را برای اعزام شما در نظر خواهد گرفت، در نتیجه در این مرحله حتما پس از معاینه
    از پزشک درباره نتیجه سوال کنید. 
    موارد متداولی که منجر به ثبت کمیسیون می‌شود شامل تتو، شماره چشم بالا، پای پرانتزی و موارد مشابه است. 

    معمولا مشکل خاصی برای افراد رخ نمی‌دهد ولی دانستن این نکات بهتر از ندانستن آن‌هاست. بهتر است در صورتی که به موردی شک دارید آن را 
    در گروه نظام‌وظیفه مطرح کنید.
    \par}
    
    \item مراجعه مجدد به پلیس +۱۰ جهت دریافت پرینت برگه سبز.
    \par\smallskip
    {\small
    دقت در صحت تاریخ اعزام درج‌شده توسط اپراتور الزامی است. پیش از دریافت برگه سبز مهر شده، یک نسخه اولیه از آن را امضا 
    خواهید کرد، این نسخه را به دقت بررسی نمایید و از درست بودن تاریخ اطمینان حاصل کنید. 
    \par}
\end{enumerate}


\subsection{معافیت تحصیلی خارج}
تکمیل شود.